\chapter{Occupation and Modular Spectra in the Hilbert--P\'olya Problem}
\label{v4-arith:w5-b:hp-operator}

The prime character uses the nonnegative occupation energies
$\log m$. The modular operator of a state acts on observables,
so its energies are differences of occupation energies. We
construct the comparison between these two operators before
asking for a Hilbert--P\'olya realization. Their relation involves
a closed creation subspace and a factor of two in the convention
$N_{\mathrm{BC}}=-\tfrac12\log\Delta_1$. It is not an
identification of the two full Hilbert spaces.

\section{The Tensor Shift State}
\label{v4-arith:w5-b:bc-setup}

\begin{definition}[tensor shift algebra and state]
\label{v4-arith:w5-b:def:bc-algebra}
Let $\mathcal T_p$ be the represented prime shift algebra of
\Cref{v4-arith:thermal:prop:normal-form}. For finite $F$ put
$\mathcal T_F=\bigotimes_{p\in F}\mathcal T_p$, and form the
algebraic union $\mathcal T_{\mathrm{fin}}$ under
$a\mapsto a\otimes1$. Its time evolution is
\begin{equation}
 \alpha_t(S_p)=p^{it}S_p,\qquad
 \alpha_t(R_p)=p^{-it}R_p.
 \label{v4-arith:w5-b:eq:bc-time}
\end{equation}
At $\beta=1$, the normalized product state on each finite factor
is $\Phi_{F,1}=\bigotimes_{p\in F}\Phi_{p^{-1}}$.
\end{definition}

\begin{lemma}[compatible normalized state]
\label{v4-arith:w5-b:lem:kms-1}
\ClaimStatusProvedHere. These functionals define a positive
normalized KMS functional $\Phi_1$ on $\mathcal T_{\mathrm{fin}}$.
Its values on the prime creation and division operators satisfy
$\Phi_1(\mu_n^*\mu_m)=\delta_{nm}$ and
$\Phi_1(\mu_n\mu_m^*)=\delta_{nm}m^{-1}$.
\end{lemma}

\begin{proof}
A finite collection of observables lies in one $\mathcal T_F$.
There positivity, normalization and the KMS identity follow by
finite tensoring from \Cref{v4-arith:thermal:prop:thermal-values}.
The state takes value one on every new unit, so the values agree
under inclusions. The normal-word values in each prime factor
give the two formulas.
\end{proof}

\begin{remark}[arithmetic phases and the convergence boundary]
\label{v4-arith:w5-b:rem:wave4-ref}
\label{v4-arith:w5-b:rem:partition-at-one}
The algebra here is the tensor shift core. The presented local
phase algebras and their represented quotients are distinguished
in \Cref{v4-arith:kms-gns:cycle1:state}; the genuine arithmetic
phase does not factor through the naive singleton inclusions by
\Cref{v4-arith:thermal:prop:two-prime-phase}. Thus this product
state is not a construction of the native adelic phase system.
Nor is it a normalized global density at $\beta=1$ on the
occupation space: its putative partition trace is the divergent
series $\sum_m m^{-1}$. A finite part of the scalar pole at one
is not a normalization of that divergent operator trace.
\end{remark}

\section{The Full Tensor GNS Spectrum}
\label{v4-arith:w5-b:gns-modular}

\begin{definition}[product GNS space]
\label{v4-arith:w5-b:def:gns}
Let $\mathcal H_1$ be the completion of the algebraic union of
finite-prime GNS spaces, using the isometries
$X\mapsto X\otimes\varrho_{\mathrm{new}}^{1/2}$.
The reference vector is the tensor of the normalized local
vectors $\Omega_p=\varrho_{p^{-1}}^{1/2}$, and the algebra acts
by left multiplication. Equivalently this is the GNS space of
$\Phi_1$, since its dense cyclic vectors have exactly those
finite-factor inner products.
\end{definition}

\begin{lemma}[the Hilbert tensor carrier]
\label{v4-arith:w5-b:lem:gns-restricted-tensor}
\ClaimStatusProvedHere. There is a unitary identification
\[
 \mathcal H_1
   =\bigotimes_{p}^{\Omega_p}\mathcal S_2(\mathscr H_p).
\]
The superscript specifies the incomplete Hilbert tensor product:
elementary tensors have factor $\Omega_p$ outside a finite set,
and their inner product is the finite product of local inner
products. Every $\Omega_p$ has norm one.
\end{lemma}

\begin{proof}
\Cref{v4-arith:thermal:gns-carrier} identifies each local GNS
factor with its Hilbert--Schmidt space. The product state inner
product on a finite set is the product inner product. The unit
observable maps to $\Omega_p$ in each new factor, so these
identifications agree with the isometric inclusions. Completing
their common dense union proves the assertion.
Finite tensors of right matrix multiplications are bounded, commute
with the left algebra, and send the reference vector to a dense
subspace. They still commute with its strong operator closure.
If an operator \(a\) in that closure kills \(\Omega_1\), it therefore
kills this dense right-cyclic subspace and is zero. Thus \(\Omega_1\)
is separating as well as cyclic for the represented von Neumann
algebra, as required for its modular operator.
\end{proof}

Choose an orthonormal basis of the diagonal Hilbert--Schmidt
subspace at each prime which contains $\Omega_p$. Such a basis
is obtained by applying Gram--Schmidt to
$\Omega_p,E_{00},E_{11},\ldots$, omitting zero remainders.
The coefficients can be real. Together with $E_{ij}$ for
$i\ne j$, this gives a local orthonormal basis. Finite tensors
of its elements, with $\Omega_p$ elsewhere, form a countable
orthonormal basis of $\mathcal H_1$. Each diagonal basis vector
has modular eigenvalue one, and each $E_{ij}$ has eigenvalue
$p^{-(i-j)}$.

\begin{lemma}[product modular operator]
\label{v4-arith:w5-b:lem:modular}
\ClaimStatusProvedHere. The positive diagonal operator determined
on this tensor basis by
\begin{equation}
 \Delta_1E_{\mathbf i,\mathbf j}
    =\prod_p p^{-(i_p-j_p)}E_{\mathbf i,\mathbf j},
 \qquad \Delta_1^{it}[x]=[\alpha_{-t}(x)]
 \label{v4-arith:w5-b:eq:Delta-basis}
\end{equation}
is the modular positive factor for the product GNS state.
A diagonal factor in the chosen basis contributes exponent zero.
The domain of $\Delta_1$ consists of expansions for which the
sum of squared coefficients times the squared displayed
eigenvalues is finite. Its conjugation $J_1$ is the continuous
extension of the finite tensor conjugations.
\end{lemma}

\begin{proof}
The local formulas are \Cref{v4-arith:thermal:modular}. Finite
products of the local eigenvalues give the formula. Positivity
and self-adjointness hold on the maximal weighted diagonal
domain. Every tensor basis vector belongs to the algebraic GNS
cyclic domain: a local off-diagonal matrix unit is a nonzero
multiple of $[E_{ij}]$, and every chosen diagonal basis vector
is a finite linear combination of $\Omega_p$ and $[E_{ii}]$.
Thus finite spans of the tensor eigenbasis lie in that domain
and form a core for $\Delta_1^{1/2}$. On the entire algebraic
cyclic domain, the finite-prime calculation gives
$[x]\mapsto[x^*]=J_1\Delta_1^{1/2}[x]$. Taking the closure
proves the polar decomposition. The group formula follows on
homogeneous cyclic vectors and extends by continuity.
\end{proof}

\begin{definition}[the specified log-modular normalization]
\label{v4-arith:w5-b:def:N-BC}
Keep the normalization
\[
 N_{\mathrm{BC}}=-\tfrac12\log\Delta_1.
\]
It is defined on the maximal diagonal domain of this real
logarithm. In one prime factor,
\begin{equation}
 N_{\mathrm{BC}}E_{ij}=\tfrac12(i-j)\log p\,E_{ij}.
 \label{v4-arith:w5-b:eq:N-BC}
\end{equation}
\end{definition}

\begin{theorem}[spectrum of the full log-modular operator]
\label{v4-arith:w5-b:thm:N-BC-self-adjoint}
\ClaimStatusProvedHere. The finite tensor eigenbasis is a core
for $N_{\mathrm{BC}}$, whose closure is self-adjoint. Its
set of eigenvalues is
\[
 \left\{\tfrac12\log(a/b):a,b\in\mathbb N_{\geq1}\right\}.
\]
Each eigenvalue has countably infinite multiplicity. The spectrum
as a closed subset of $\mathbb R$ is all of $\mathbb R$,
although every spectral measure is atomic. Moreover
$J_1N_{\mathrm{BC}}J_1=-N_{\mathrm{BC}}$.
In particular the full log-modular operator is not nonnegative.
\end{theorem}

\begin{proof}
The logarithm of a finite product of prime powers is the
logarithm of a positive rational number; conversely every such
rational is obtained from finite positive and negative occupation
differences. This gives exactly the stated eigenvalues. Keeping
a fixed difference and varying the common diagonal shift of
$i,j$ gives infinitely many orthogonal eigenvectors; separability
makes the multiplicity countable. The positive rational numbers
are dense in the positive real numbers, and logarithm is
continuous, so the eigenvalue set is dense in $\mathbb R$.
For a diagonal self-adjoint operator, every limit of eigenvalues
lies in the spectrum: unit eigenvectors at approaching eigenvalues
are approximate null vectors for the shifted operator. A point
at positive distance from the eigenvalues has a bounded diagonal
resolvent. These two facts identify the spectrum with their
closure. Self-adjointness and the core assertion follow by
truncating the weighted square sum of coefficients.
Finally $J_1$ reverses each occupation difference. Already in a
single prime, $E_{10}$ and $E_{01}$ have opposite nonzero energies.
\end{proof}

\section{The Creation Subspace and the Occupation Hamiltonian}
\label{v4-arith:w5-b:primary-heisenberg}

Let $\mathcal H_{\mathrm{occ}}=\ell^2(\mathbb N_{\geq1})$
with orthonormal basis $v_m$. Define
\[
 N_{\mathrm{occ}}v_m=\log m\,v_m,\qquad
 \operatorname{Dom}N_{\mathrm{occ}}
   =\left\{\sum_m c_mv_m:\sum_m(\log m)^2|c_m|^2<\infty\right\}.
\]
The same diagonal adjoint test proves self-adjointness. Its
spectrum is $\{\log m:m\geq1\}$, with multiplicity one, and
it is nonnegative.

\begin{theorem}[occupation-to-GNS comparison]
\label{v4-arith:w5-b:occupation-comparison}
\ClaimStatusProvedHere. The map
\[
 W:\mathcal H_{\mathrm{occ}}\longrightarrow\mathcal H_1,
 \qquad Wv_m=[\mu_m]
\]
is an isometry onto the closed creation subspace. It intertwines
prime creation and satisfies
\[
 WN_{\mathrm{occ}}=2N_{\mathrm{BC}}W
\]
with equality of the maximal domains on that subspace. The
creation subspace is reducing for $N_{\mathrm{BC}}$ but is not
invariant under $J_1$. The full GNS algebra of observables does
not preserve it.
\end{theorem}

\begin{proof}
The state formula gives
$\langle[\mu_n],[\mu_m]\rangle=\delta_{nm}$, so $W$ is an
isometry on the algebraic basis and extends to its completion.
Left creation sends $[\mu_m]$ to $[\mu_{pm}]$, proving its
intertwining relation. Modular covariance gives
$\Delta_1[\mu_m]=m^{-1}[\mu_m]$, hence the logarithmic identity.
The closed span of these eigenvectors is reducing. Equality of
the maximal domains follows from the same weighted coefficient
sum on both sides.

For $p>1$, $J_1[\mu_p]=p^{1/2}[\mu_p^*]$. The division vector
is nonzero and orthogonal to every creation vector, because its
inner products involve a nonconstant pure creation word.
It therefore lies outside the creation subspace. Also left
$\mu_p^*$ sends $\Omega_1$ to this division vector, so the
full observable algebra does not preserve that subspace.
\end{proof}

\begin{lemma}[the weighted zero-mode embedding]
\label{v4-arith:w5-b:lem:zero-mode-embeds-gns}
\ClaimStatusProvedHere. The weighted lattice with basis
$\delta_m$ of squared norm $m^{-1}$ has the isometric embedding
\[
 \iota_0\delta_m=m^{-1/2}[\mu_m].
\]
It lies in the creation subspace and has the same occupation
energies there under $2N_{\mathrm{BC}}$.
\end{lemma}
\begin{proof}
Rescale the orthonormal family of the preceding theorem.
The scalar rescaling does not change its eigenvalues.
\end{proof}

\begin{remark}[the exact carrier distinction]
\label{v4-arith:w5-b:rem:N-BC-not}
\label{v4-arith:w5-b:rem:embedding-domain}
The original positive log-prime target is realized by
$N_{\mathrm{occ}}$ and by $2N_{\mathrm{BC}}$ on the creation
subspace. It is not the full modular spectrum. No unitary
intertwiner can identify the full $N_{\mathrm{BC}}$ with the
nonnegative occupation operator, since the former has negative
eigenvalues. Nor is the creation subspace a substitute for the
full GNS space when division, modular conjugation, or genuine
arithmetic phases are required.
\end{remark}

\begin{definition}[an oscillator enlargement and its lift]
\label{v4-arith:w5-b:def:fock}
\label{v4-arith:w5-b:def:bc-gns-heisenberg-lift}
An oscillator enlargement is a specified Hilbert completion of
finite-support number states $|\mathbf n\rangle\otimes\delta_m$,
with creation and annihilation operators satisfying the chosen
Heisenberg relations, a vacuum and a nonnegative diagonal energy
$L_0^{\mathrm{osc}}$. Its zero-mode energy is $N_{\mathrm{occ}}$.
A GNS lift of this enlargement is an isometry of its number-state
core into $\mathcal H_1$ extending $\iota_0$, intertwining the
specified operators on their domains and preserving the given
inner products. The notation
$\mathcal F^{\mathrm{Ar}}_{\mathrm{Heis},\mathrm{prim}}$ denotes
such a chosen core when one is supplied.
\end{definition}

\begin{lemma}[completion of an oscillator lift]
\label{v4-arith:w5-b:lem:fock-embeds-gns}
\ClaimStatusProvedHereConditional. An oscillator lift extends
uniquely to an isometry of its Hilbert completion onto the closed
span of its image. It supplies no identification of that image
with a native arithmetic Hochschild or automorphic carrier.
\end{lemma}
\begin{proof}
An isometry sends Cauchy sequences to Cauchy sequences and
preserves the norm of their limits. This gives the unique
extension and identifies its closed image. The assertion uses
only the specified lift maps, not a classification of additional
GNS coordinates as automorphic sectors.
\end{proof}

\section{The Primary-Heisenberg Twist and the Determinant-Trace Operator}
\label{v4-arith:w5-b:twist}

The occupation operator \(N_{\mathrm{occ}}\) has spectrum
\(\{\log m\}_{m\in\mathbb{N}^{\times}}\). The Hilbert--P\'olya
condition asks for a self-adjoint operator \(A\) together with
the determinant condition forcing the divisor of
\(\tfrac12\mathrm{id}+iA\) to be the zero divisor of
\(\xi_{\mathbb R}\). The \emph{primary-Heisenberg twist} names the
data required to compare the first spectrum with this determinant
divisor and its Weil--Connes trace. Its existence with the prescribed
determinant-trace property is H-P\(^{\mathrm{Ar}}\).

\subsection{Statement of the Twist}
\label{v4-arith:w5-b:twist:statement}

\begin{definition}[graded determinant-trace comparison data]
\label{v4-arith:w5-b:def:twist}
Let $\mathcal F$ be the Hilbert completion of a specified
oscillator enlargement, with its dense number-state core and
self-adjoint grading operator
$G_{\mathcal F}=L_0^{\mathrm{osc}}+N_{\mathrm{occ}}$.
A graded determinant-trace comparison consists of a Hilbert
space $\mathcal H_\xi$, a unitary $U:\mathcal F\to\mathcal H_\xi$,
and a self-adjoint operator $A$ on $\mathcal H_\xi$ with a
specified invariant core $\mathcal D_\xi$. The transported
grading operator is
\begin{equation}
 G_\xi=UG_{\mathcal F}U^*,\qquad
 G_\xi Uv=UG_{\mathcal F}v
       \quad(v\in\operatorname{Dom}G_{\mathcal F}),
 \label{v4-arith:w5-b:eq:twist-intertwine}
\end{equation}
with $\operatorname{Dom}G_\xi=U\operatorname{Dom}G_{\mathcal F}$.
If $\mathcal H_\xi$ is to be a native Petersson completion, its
pairing realization maps are also required. The operator $A$
is a separate operator from $G_\xi$; no equality or intertwining
between them is included in the definition. Write
\begin{align}
 H_{\mathrm{HP}}&=A,
 \label{v4-arith:w5-b:eq:H-HP-def}\\
 \Theta_\xi&=\tfrac12\mathrm{id}+iA.
 \label{v4-arith:w5-b:eq:Theta-xi-def}
\end{align}
The determinant and trace requirements on $A$ are additional
hypotheses below.
\end{definition}

\begin{proposition}[the literal energy-intertwining obstruction]
\label{v4-arith:w5-b:rem:what-iv-says}
\label{v4-arith:w5-b:vacuum-obstruction}
Suppose the source grading has a nonzero vacuum $v_1$ with
$G_{\mathcal F}v_1=0$. Requiring
$AUv=UG_{\mathcal F}v$ on a core containing $v_1$ is incompatible
with the requirement that the eigenvalue divisor of
$\tfrac12+iA$ equal the nontrivial zero divisor of $\xi_{\mathbb R}$.
\end{proposition}

\begin{proof}
The proposed identity would give $AUv_1=0$, and $Uv_1\ne0$
because $U$ is an isometry. Hence $1/2$ would be an eigenvalue
of $\tfrac12+iA$ and would have to be a zero of $\xi_{\mathbb R}$.
It is not. Indeed the alternating series
\[
 \eta(1/2)=\sum_{k\geq1}
       \bigl((2k-1)^{-1/2}-(2k)^{-1/2}\bigr)
\]
converges and is positive: every term is positive, and the mean
value estimate bounds the terms by a constant times $k^{-3/2}$.
The identity
$\zeta(s)=\eta(s)/(1-2^{1-s})$ follows first by separating even
terms for $\operatorname{Re}s>1$ and continues near $s=1/2$
by uniform Dirichlet convergence of the alternating series.
Its denominator at $1/2$ is nonzero, so $\zeta(1/2)\ne0$.
The gamma factor at $1/4$ is positive by its defining integral,
and removing the Tate poles multiplies by the nonzero factor
$\tfrac12s(s-1)$ at $s=1/2$. Thus $\xi_{\mathbb R}(1/2)\ne0$.
\end{proof}

The grading comparison \eqref{v4-arith:w5-b:eq:twist-intertwine}
and the literal $A$-intertwining condition are therefore different
targets. The former is explicitly constructed once $U$ is given.
The latter is ruled out on a vacuum-preserving domain. Neither
statement constructs, or excludes in general, a self-adjoint
operator with the required zeta-zero determinant and trace.

\subsection{Determinant-Trace Requirement: Hadamard Inversion}
\label{v4-arith:w5-b:twist:hadamard-inversion}

The condition on \(A\) is the Hilbert--P\'olya determinant-trace
condition for \(\xi_{\mathbb R}\). Write a non-trivial zero as
\(\rho=\beta_{\rho}+i\gamma_{\rho}\). The self-adjointness of \(A\)
puts the normal operator \(\Theta_{\xi}=1/2+iA\) on the critical line.
H-P\(^{\mathrm{Ar}}\) asks that its regularised determinant divisor be
the original zero divisor and that its Weil--Connes trace be the zero
distribution. Equivalently, the determinant identity forces
\(\beta_{\rho}=1/2\) for every zero; after this zero-placement the
trace of \(A\) is supported at the ordinates \(\gamma_{\rho}\). This is
\emph{not} the spectrum of \(N_{\mathrm{occ}}\), which is
\(\{\log m\}_{m\in\mathbb{N}^{\times}}\). The prime and zero sides are
related by the Mellin transform of \(\xi\) through its Hadamard
factorisation:
\begin{equation}
\label{v4-arith:w5-b:eq:hadamard}
\xi_{\mathbb R}(s)\;=\;\frac{1}{2}\,e^{B s}\,\prod_{\rho}\Bigl(1-\frac{s}{\rho}\Bigr)e^{s/\rho}
\end{equation}
(Riemann 1859; Hadamard 1893), where the product is over non-trivial
zeros and $B=\xi'(0)/\xi(0)=-1-\gamma/2+\tfrac12\log(4\pi)$,
as computed in Proposition~\ref{v4-arith:pm:constant}.
For $\Re s>1$, the primes appear through the absolutely convergent
Dirichlet series in the \emph{logarithmic derivative}
\begin{equation}
\label{v4-arith:w5-b:eq:log-deriv-xi}
-\frac{\xi_{\mathbb R}'(s)}{\xi_{\mathbb R}(s)}\;=\;-\frac{1}{s}-\frac{1}{s-1}+\frac{1}{2}\log\pi-\frac{1}{2}\Gamma'/\Gamma(s/2)+\sum_{n\geq1}\frac{\Lambda(n)}{n^{s}}
\end{equation}
(von Mangoldt 1895). Outside $\Re s>1$ the prime series is not
silently retained: the continued logarithmic derivative is a
meromorphic function, while the Weil formula uses the compact-test
pairing and exponential weight of Theorem~\ref{v4-arith:pm:weil}. The twist \(A\) is the operator-level analogue of
\emph{inverting this duality}: from the log-prime spectrum of
\(N_{\mathrm{occ}}\), construct a normal zero-divisor operator whose
determinant is governed by the Hadamard expansion and whose trace is
the Weil--Connes zero distribution. The Hadamard expansion is a
statement about an entire function. The assertion that a self-adjoint
operator carries this determinant divisor and trace distribution is
additional data, the content of H-P\(^{\mathrm{Ar}}\).

\subsection{The Hilbert--P\'olya Hypothesis and Its Determinant-Trace Consequences}
\label{v4-arith:w5-b:twist:consequences}

\begin{hypothesis}[H-P\(^{\mathrm{Ar}}\)-twist]
\label{v4-arith:w5-b:hyp:twist-existence}
There exist graded determinant-trace comparison data \((\mathcal{H}_{\xi},U,A)\) in
the sense of \Cref{v4-arith:w5-b:def:twist} with the additional
zero-placement spectral property
\begin{equation}
\label{v4-arith:w5-b:eq:spec-A}
\mathrm{Spec}(A)\;=\;\{\gamma_{\rho}\mid\rho=\beta_{\rho}+i\gamma_{\rho}
\text{ is a non-trivial zero of }\xi_{\mathbb R}\},
\end{equation}
with multiplicities matching zero multiplicities, together with the
determinant condition
\(\mathrm{div}\,\det_{\infty}(s\mathrm{id}-(\tfrac12\mathrm{id}+iA))
=\mathrm{div}\,\xi_{\mathbb R}\), and the
\(\mathbb{Z}/2\)-symmetry \(\gamma_{\rho}\mapsto-\gamma_{\rho}\) acting on
the spectrum as the arithmetic \(S\)-modular involution.
\end{hypothesis}

\begin{remark}[relation to H-P\(^{\mathrm{Ar}}\)]
\label{v4-arith:w5-b:rem:equals-HPAr}
\Cref{v4-arith:w5-b:hyp:twist-existence} is the functional-analytic
content of \Cref{v4-arith:rh:hyp:HPAr}: an operator \(A\) with the
specified determinant divisor and trace distribution is exactly
\(H_{\mathrm{HP}}\) of the reformulation chapter, and
\(\Theta_{\xi}=1/2+iA\) is the arithmetic zero-divisor operator.
The arithmetic realization asks in addition for the specified occupation
grading, native maps and arithmetic trace comparison. The grading
intertwiner \(U\) in \Cref{v4-arith:w5-b:def:twist} does not intertwine
\(A\) with \(N_{\mathrm{occ}}\). Imposing the latter identity on the
vacuum contradicts \Cref{v4-arith:w5-b:vacuum-obstruction}.
\end{remark}

\section{Determinant Identity: Hadamard Factorisation}
\label{v4-arith:w5-b:determinant}

\subsection{The normalized genus-one determinant}
\label{v4-arith:sc:determinant}

Let $Z$ contain each zero $\rho$ of $\xi$ with its multiplicity.
On $K_Z=\ell^2(Z)$, define
\[
 (\Theta_Zu)_\rho=\rho u_\rho,\qquad
 \operatorname{Dom}\Theta_Z=\{u:\sum_\rho|\rho u_\rho|^2<\infty\}.
\]
This is a closed normal operator. Its inverse is bounded, compact,
and Hilbert--Schmidt: $\xi(0)=1/2$, the zeros are discrete, and
$\sum_\rho|\rho|^{-2}<\infty$.

\begin{theorem}[complete determinant normalization]
\label{v4-arith:sc:normalized-det}
\label{v4-arith:w5-b:prop:det-formal}
Define the second regularized Fredholm determinant of the diagonal
Hilbert--Schmidt operator $-s\Theta_Z^{-1}$ by
\[
 D_2(s)=\det{}_2(1-s\Theta_Z^{-1})
       :=\prod_{\rho\in Z}(1-s/\rho)e^{s/\rho}.
\]
With $B=-1-\gamma/2+\tfrac12\log(4\pi)$, set
\begin{equation}
 \Delta_Z(s):=\tfrac12e^{Bs}D_2(s).
 \label{v4-arith:sc:normalized-det-formula}
\end{equation}
This entire function equals $\xi(s)$.
It is the unique entire function of order at most one with this
zero divisor, $\Delta_Z(0)=1/2$, and $\Delta_Z'(0)/\Delta_Z(0)=B$.
Off the zero divisor it satisfies
\[
 (\log\Delta_Z)''(s)
       =-\operatorname{Tr}(\Theta_Z-s)^{-2}.
\]
The trace is an ordinary absolutely convergent trace.
Unitary conjugacy of $\Theta_Z$, including its domain, preserves
this determinant prescription.
\end{theorem}
\begin{proof}
For $|w|\le1/2$, $\log(1-w)+w=O(|w|^2)$.
The inverse-square sum proves normal convergence of the product
on every compact set, independently of the enumeration of zeros.
Its value and derivative at zero are $1$ and $0$.
Proposition~\ref{v4-arith:pm:constant} gives the product for $\xi$
with exactly the factor $e^{Bs}/2$. This proves the identity.
Hadamard factorization writes every other entire function of order
at most one with this divisor as $e^{a+bs}D_2(s)$.
The two initial conditions give $e^a=1/2$ and $b=B$.
Twice differentiating the normally convergent logarithmic product
gives $-\sum_\rho(\rho-s)^{-2}$.
The inverse-square estimate proves absolute trace convergence.
Unitary conjugacy preserves the diagonal eigenvalues and their
multiplicities, hence every factor in the prescribed product.
\end{proof}

\begin{proposition}[the divisor leaves an exponential ambiguity]
\label{v4-arith:sc:det-ambiguity}
For every $c\in\mathbb C$, the function $e^{cs}\xi(s)$ has the same
divisor, the same value at zero, and order at most one.
Its logarithmic derivative at zero is $B+c$.
Thus a divisor condition and a value normalization do not imply
the determinant identity. Both normalizations in
\eqref{v4-arith:sc:normalized-det-formula} are necessary.
\end{proposition}
\begin{proof}
The exponential has no zero, has value one at zero, and contributes
$c$ to the logarithmic derivative. Multiplication preserves the
stated order bound.
\end{proof}

The prescription above uses no complex powers of $s-\Theta_Z$.
A determinant defined by such powers requires a spectral cut and
meromorphic continuation of its trace. Agreement with $\Delta_Z$
requires the same second logarithmic derivative and the two stated
initial conditions. These conditions determine equality locally
near zero by integration, and globally by analytic continuation.

\begin{remark}[operator comparison and determinant normalization]
\label{v4-arith:w5-b:rem:det-does-not-buy}
For a normal operator with a complete orthonormal eigenbasis indexed
by the same zeros, use the prescription
$\det_{\infty}(s-\Theta)=\tfrac12e^{Bs}\det_2(1-s\Theta^{-1})$.
The unitary eigenbasis map proves equality with $\Delta_Z$.
This fixes the determinant notation in the determinant-trace
hypothesis. The construction of such an operator from an arithmetic
primary sector requires an additional domain-preserving comparison.
The diagonal operator $\Theta_Z$ itself uses the zero divisor as input.
\end{remark}

\begin{remark}[comparison with Berry--Keating]
\label{v4-arith:w5-b:rem:berry-keating-contrast}
Berry--Keating (SIAM Review 41, 1999) propose \(H=xp+1/2\) on
\(L^{2}(\mathbb{R}_{>0},dx/x)\) as the candidate \(H_{\mathrm{HP}}\).
Their proposal reproduces the leading-order Riemann--von Mangoldt
zero density via Bohr--Sommerfeld quantisation; the Planck-cell
regularisation, however, requires a self-adjoint extension not supplied
by the model. The present twist has the same logical form: the occupation
log-prime spectrum \(\{\log m\}_{m\in\mathbb{N}^{\times}}\) is
classical and self-adjoint, and the lift to the zero determinant
divisor via the Hadamard-dual twist is the content of
H-P\(^{\mathrm{Ar}}\). See
\Cref{v4-arith:hpvbconv:rem:berry-keating} for the HPAr-converse analysis.
\end{remark}

\section{Trace Distribution: Weil--Connes Explicit Formula}
\label{v4-arith:w5-b:trace}

\begin{proposition}[the trace on the compact Fourier test class]
\label{v4-arith:w5-b:prop:weil-trace}
\ClaimStatusProvedHereConditional. Suppose the spectral realization
of Hypothesis~\ref{v4-arith:w5-b:hyp:twist-existence} has been
constructed: $H_{\mathrm{HP}}$ is self-adjoint with a complete
orthonormal eigenbasis indexed by the actual xi zeros, including
multiplicity, and $\rho=1/2+i\gamma_\rho$ is the eigenvalue of its
shifted operator. For $f\in C_c^\infty(\mathbb R)$ put
$h(z)=\int f(t)e^{izt}dt$. Then $h(H_{\mathrm{HP}})$ is trace-class
and
\begin{align}
\label{v4-arith:w5-b:eq:trace-identity}
 \operatorname{Tr}h(H_{\mathrm{HP}})
 &=\sum_\rho h(\gamma_\rho)
   =h(i/2)+h(-i/2)+\mathcal A_\infty(h)\notag\\
 &\quad-\sum_{m\geq2}\frac{\Lambda(m)}{\sqrt m}
                         (f(\log m)+f(-\log m)),
\end{align}
where
\[
 \mathcal A_\infty(h)=\frac1{2\pi}\int_{\mathbb R}
 h(u)(\Re\psi_\Gamma(1/4+iu/2)-\log\pi)\,du.
\]
If a specified regularized trace agrees with the ordinary trace
on trace-class operators, it has the same value here.
Extension to a larger Paley--Wiener class requires the domain and
continuity of that trace functional on the larger class.
\end{proposition}
\begin{proof}
Repeated integration by parts makes $h$ rapidly decreasing on the
real line. The zero count of Lemma~\ref{v4-arith:pm:zero-count}
therefore gives $\sum_\rho|h(\gamma_\rho)|<\infty$.
Functional calculus on the stipulated complete eigenbasis proves
the trace-class assertion and the first equality.
The actual-divisor hypothesis identifies $\gamma_\rho$ with
$-i(\rho-1/2)$, so Theorem~\ref{v4-arith:pm:weil} proves the rest.
The compatibility hypothesis gives the regularized-trace clause;
it is not a consequence of an isometric embedding or of the scalar
determinant normalization. For even $f$ the last sum reduces to
$2\mu_{1/2}(f)$, with precisely the stated inverse Fourier convention.
\end{proof}

\begin{corollary}[positivity on the complete Fourier square class]
\label{v4-arith:w5-b:cor:weil-positive}
\ClaimStatusProvedHereConditional. Under the preceding realization,
let $g\in C_c^\infty(\mathbb R)$,
$\phi(z)=\int g(t)e^{izt}dt$, and
$h(z)=\phi(z)\overline{\phi(\overline z)}$.
Then $h$ is the Fourier transform of $g*g^*$ and
\[
 \operatorname{Tr}h(H_{\mathrm{HP}})
                 =\sum_\rho|\phi(\gamma_\rho)|^2\geq0.
\]
Thus the actual Weil functional is positive on this entire square
class. The realization itself implies RH, and hence classical Li
positivity. Identifying a Petersson-positive cone with this full
test class requires a separately specified map and its density and
continuity properties.
\end{corollary}
\begin{proof}
The convolution and involution identities of
Definition~\ref{v4-arith:mc:test-algebra} give the stated entire
function $h$. The preceding trace formula now has nonnegative
eigenvalues and absolutely convergent sum. The real spectrum of
$H_{\mathrm{HP}}$ places every actual zero on $\Re s=1/2$;
Li positivity follows from the classical Li criterion
in Theorem~\ref{v4-arith:rh:thm:Li-criterion}.
An identity known only on an unspecified restricted cone does not
supply the comparison with all compact Fourier squares.
\end{proof}

\section{Operator and Arithmetic Comparison Boundary}
\label{v4-arith:w5-b:form}
\label{v4-arith:w5-b:form:layer-1}
\label{v4-arith:w5-b:form:layer-2}
\label{v4-arith:w5-b:form:layer-3}
\label{v4-arith:w5-b:form:summary}

The tensor shift state at $\beta=1$ has a well-defined GNS
space and modular operator. Its full log-modular spectrum is
$\mathbb R$, with an atomic spectral resolution on the dense
set $\tfrac12\log\mathbb Q_{>0}$. The occupation operator has
the discrete nonnegative spectrum $\log\mathbb N_{\geq1}$.
The map $W$ of \Cref{v4-arith:w5-b:occupation-comparison}
identifies the latter with $2N_{\mathrm{BC}}$ on the creation
subspace. Division and modular conjugation require the larger
GNS carrier. Genuine arithmetic phases require their own
compatible algebra before an adelic comparison can be made.

\begin{proposition}[the remaining determinant-trace construction]
\label{v4-arith:w5-b:thm:precise-obstruction}
The local GNS and occupation constructions supply neither an
operator $A$ with the zero-divisor property of
\Cref{v4-arith:w5-b:hyp:twist-existence} nor a native pairing
map into its Hilbert space. The requested additional data are
that self-adjoint operator and core, the specified determinant
and trace regularizations, and the native realization maps.
The separate grading comparison is
\eqref{v4-arith:w5-b:eq:twist-intertwine}; imposing it on $A$
itself is excluded by \Cref{v4-arith:w5-b:vacuum-obstruction}.
\end{proposition}

\begin{proof}
The constructions give the displayed diagonal operators and the
isometry $W$. Their spectra, carriers and actions are computed
above. The occupation operator retains its vacuum and the full
modular operator has both signs and infinite eigenspace
multiplicity. Neither has the asserted zeta-zero divisor.
Transport of the grading only constructs $G_\xi$. It imposes
no spectral condition on the separate operator $A$. The proposed
stronger energy intertwining is the vacuum contradiction already
proved. Thus the additional hypotheses remain the exact
unconstructed comparison data, not consequences of the local
state normalization.
\end{proof}

\begin{remark}[prime weights and native realization]
\label{v4-arith:w5-b:rem:W-honest}
The weight $k\log p$ is the occupation energy of $p^k$.
It is obtained from $2N_{\mathrm{BC}}$ only after the stated
creation-subspace restriction at $\beta=1$; at general inverse
temperature $\beta$ the factor is $2/\beta$.
A prescribed scalar partition determines neither the native
pairing maps nor the full algebra of phases. In particular a
local KMS normalization cannot establish uniqueness of an
arithmetic Hilbert--P\'olya representative.
\end{remark}

\begin{remark}[scope of the Hilbert--P\'olya implication]
\label{v4-arith:w5-b:rem:honest-conditional}
If a self-adjoint $A$ has the required determinant divisor and
trace, then $\tfrac12+iA$ is supported on the critical line and
the assumed divisor comparison transfers this support to the
zeros. Constructing such an $A$ remains the Hilbert--P\'olya
problem in the specified carrier. The present normalization
calculations do not decide that existence problem or the
positivity of a distinct native arithmetic operator. Comparisons
with absorption traces or motivic realizations additionally
require maps to those particular carriers.
\end{remark}

\section{Bibliography and Source List}
\label{v4-arith:w5-b:biblio}

The argument draws on four disjoint source lists
here with publisher/arXiv identifiers for reproducibility.

\paragraph{Operator-algebra and KMS theory.}
\begin{description}
\item[Bost--Connes.] J.-B.~Bost and A.~Connes, \emph{Hecke algebras,
type~III factors and phase transitions with spontaneous symmetry
breaking in number theory}, Selecta Math.~1 (1995), 411--457.
\item[Connes.] A.~Connes, \emph{Trace formula in noncommutative
geometry and the zeros of the Riemann zeta function}, Selecta
Math.~5 (1999), 29--106; arXiv:math/9811068.
\item[Connes--Marcolli.] A.~Connes and M.~Marcolli, \emph{Noncommutative
geometry, quantum fields, and motives}, AMS Colloquium Publications
55 (2008). Ch.~3--4 for KMS states and rigidity; Ch.~5 for the
adele class space.
\item[Takesaki.] M.~Takesaki, \emph{Tomita's theory of modular
Hilbert algebras and its applications}, Lecture Notes in Mathematics
128 (Springer 1970); \emph{Theory of Operator Algebras~I}, (Springer
1979), Ch.~IV \S5 (GNS), Ch.~VI \S1 (Tomita--Takesaki); \emph{Theory
of Operator Algebras~II} (Springer 2003), Ch.~VIII \S1 (modular
automorphism).
\item[Meyer.] R.~Meyer, \emph{On a representation of the idele class
group related to primes and zeros of \(L\)-functions}, Duke Math.~J.
127 (2005), 519--595; arXiv:math/0311468.
\end{description}

	\paragraph{Hilbert--P\'olya candidates and determinant-trace realisations.}
\begin{description}
\item[Berry--Keating.] M.~V.~Berry and J.~P.~Keating, \emph{The
Riemann zeros and eigenvalue asymptotics}, SIAM Review 41 (1999),
236--266.
\item[Reed--Simon.] M.~Reed and B.~Simon, \emph{Methods of Modern
Mathematical Physics}, Vol.~I (Academic Press 1980),
Thm.~VIII.11 (diagonal operators); Vol.~II (Academic Press 1975),
Thm.~X.23 (Friedrichs extension).
\item[Deninger.] C.~Deninger, \emph{Motivic \(L\)-functions and
regularised determinants}, Proc.~Symp.~Pure Math.~55.1 (1994),
707--743; \emph{A note on arithmetic topology and dynamical systems},
Math.~Nachr.~205 (2001), 107--135.
\item[Ray--Singer.] D.~B.~Ray and I.~M.~Singer, \emph{\(R\)-torsion
and the Laplacian on Riemannian manifolds}, Advances in
Mathematics 7 (1971), 145--210 (source of \(\zeta\)-regularised
determinants).
\item[Voros.] A.~Voros, \emph{Spectral functions, special functions
and the Selberg zeta function}, Comm.~Math.~Phys.~110 (1987),
439--465.
\end{description}

\paragraph{Hadamard factorisation and explicit formula.}
\begin{description}
\item[Riemann.] B.~Riemann, \emph{\"Uber die Anzahl der Primzahlen
unter einer gegebenen Gr\"osse}, Monatsberichte der Berliner
Akademie, 1859.
\item[Hadamard.] J.~Hadamard, \emph{\'Etude sur les propri\'et\'es
des fonctions enti\`eres et en particulier d'une fonction consid\'er\'ee
par Riemann}, J.~Math.~Pures Appl.~9 (1893), 171--215.
\item[von Mangoldt.] H.~von Mangoldt, \emph{Zu Riemann's Abhandlung
``\"Uber die Anzahl der Primzahlen unter einer gegebenen Gr\"osse''},
J.~Reine Angew.~Math.~114 (1895), 255--305.
\item[Weil.] A.~Weil, \emph{Sur les ``formules explicites'' de la
th\'eorie des nombres premiers}, Comm.~S\'em.~Math.~Univ.~Lund (1952),
252--265.
\item[Edwards.] H.~M.~Edwards, \emph{Riemann's Zeta Function},
Academic Press 1974.
\end{description}

\paragraph{Chiral algebra and arithmetic Fock background.}
\begin{description}
\item[Beilinson--Drinfeld.] A.~Beilinson and V.~Drinfeld,
\emph{Chiral Algebras}, AMS Colloquium Publications 51 (2004).
\item[Bost--Connes~(cited above).]
\item[Vol~IV chapters.] \Cref{v4-arith:kms-gns-full},
\Cref{v4-arith:HPAr-vbconverse}, \Cref{v4-arith:hpar-det},
\Cref{v4-arith:R-analytic-comparison}, \Cref{v4-arith:rh-reformulation}.
\end{description}
